% !PRINTED 29
\weberpage{29}
In the other case we will denote the mutually distinct ones among the groups (7) by
\weq{8}{Q_\nu, \; Q'_\nu, \; Q''_\nu \; \ldots}
and from them derive a group
\weq{9}{R = (Q_\nu, \; Q'_\nu, \; Q''_\nu \; \ldots)}
which shall be defined by this, that it is the smallest group that contains each of the groups (8) as a divisor, and which we may denote as the \gesperrt{least common multiple} of the groups $Q_\nu, Q'_\nu, Q''_\nu \ldots$.

That there is one and only one such group $R$, and that every group that is divisible by $Q_\nu, Q'_\nu, Q''_\nu \ldots$ must also be divisible by $R$, is easy to see. For, first, there are finite groups, e.g. the group $P_\nu$, that contain all the groups (8) as divisors, and secondly, if two groups $R, R'$ contain all these groups, the same holds also of the intersection of $R$ and $R'$. If therefore $R$ is of the lowest possible degree, this intersection must be identical with $R$, and $R$ is thus a \gesperrt{divisor of} $R'$. If $R'$ is also of the lowest possible degree, then $R'$ is not different from $R$.

Of this group $R$ we now prove:

a) that it is a normal divisor of $P_{\nu-1}$,

b) that it is a commutative group.

When we have proved both of these, we are at our goal; for then, if there is not also a commutative normal divisor of $P_{\nu-1}$ of \gesperrt{lower degree}, we can choose $R$ itself for $Q_{\nu-1}$.

But the first is clear without further ado. For if $\gamma$ denotes any element of $P_{\nu-1}$, then
\[ \gamma^{-1} R \gamma = (\gamma^{-1} Q_\nu \gamma, \; \gamma^{-1} Q'_\nu \gamma, \; \gamma^{-1} Q''_\nu \gamma \; \ldots) \]
is the least common multiple of the groups
\[ \gamma^{-1} Q_\nu \gamma, \; \gamma^{-1} Q'_\nu \gamma, \; \gamma^{-1} Q''_\nu \gamma \; \ldots \]
But these groups, by the definition (7), are in their totality not different from the groups
\[ Q_\nu, \; Q'_\nu, \; Q''_\nu \; \ldots \]
and consequently
\[ \gamma^{-1} R \gamma = R, \]
that is, $R$ is a normal divisor of $P_{\nu-1}$.
